% Ad Familiares 6.7 --- headnote
% ad-familiares-06-07, mid-December 46 BC, Sicily

\textit{Aulus Caecina to Cicero, written from Sicily in
mid-December 46 BC (Perseus: \textit{in Sicilia med.~m.~Dec.~708
(46)}). This is the only piece of the Caecina cluster in
Caecina's own voice. For Caecina --- the antiquarian son of
Cicero's old client, the man who had circulated the
\textit{Querelae} against Caesar before the war, and who was now
in Sicilian exile waiting on Cicero's mediation --- see the
headnote to \textit{Fam.}~6.6. The letter is a reply to
\textit{Fam.}~6.5 and is sent along with a new book of Caecina's,
which his son is to put into Cicero's hand on condition that
Cicero will go over it and amend (or hold back) whatever might
do its author further harm.}

\textit{The voice is markedly Caecina's, not Cicero's: anxious,
self-aware, learned. He is exile-from-the-elite rather than
mere supplicant, and he reasons with Cicero as an equal in
cultivation about his own dangerous situation. The
antiquarian register surfaces in \S~1 in a memorable Etruscan-style
sententia, set as a tricolon of corrections: a slip of the pen
is erased, a folly is fined in reputation, ``\textit{meus error
exsilio corrigitur}'' --- my own error is corrected by exile. It
surfaces again in \S~3, in the architectural image of the
staircase whose steps have been pulled out or cut through, which
no longer carries a climber but builds the danger of collapse:
a designer's image, for the writer who could no longer write
freely about Caesar. The dramatic interior monologue of \S~4
(``This he will accept; that word is dangerous. What if I change
this?'') is the most personal moment in the cluster --- it shows
the literary man at work and afraid.}

\textit{The substantive request is in \S\S~5--6: the time for
decision has come; Cicero is not to wait on Caecina's young son,
who is too inexperienced to manage it; the whole load is
Cicero's. Caecina's careful insistence that this is not a
favour Cicero is doing but a duty he is taking up --- ``not that
you do what you are asked, but that the whole load is yours''
--- is the polite assumption of equal standing under unequal
fortune. The book mentioned at the head and foot of the letter
appears to be a new piece of Caecina's writing, and it is here,
not to the \textit{Querelae}, that the corrections refer.}
